\chapter{Конструкторская часть}

\section{Архитектура метода}

Разрабатываемый метод обеспечения дифференциальной приватности при генерации синтетических основан на Байесовских сетях и состоит из двух основных этапов:
\begin{enumerate}[label*=---]
	\item построение структуры Байесовской сети с дифференциальной приватностью;
	\item генерация синтетических данных.
\end{enumerate}

На рисунке~\ref{img:idef1} представлена функциональная модель разрабытыаемого метода.

\includeimage
{idef1} % Имя файла без расширения (файл должен быть расположен в директории inc/img/)
{f} % Обтекание (с обтеканием)
{h!} % Положение рисунка (см. wrapfigure из пакета wrapfig)
{1\textwidth} % Ширина рисунка
{Функциональная модель разрабытыаемого метода} % Подпись рисунка

\section{Предварительная обработка данных}

Тип данных является важным компонентом домена атрибутов.
Метод поддерживает четыре типа данных и позволяет пользователям явно указывать их.
Если тип данных атрибута не указан пользователем, то метод делает это автоматически.

% Для каждого атрибута сначала определяется, является ли он числовым.
% Если да, то является ли это целым или плавающим числом.
% Если атрибут не является числовым, метод пытается разобрать его как datetime. Любой атрибут, который не является ни числовым, ни datetime, считается строкой.

Тип данных атрибута ограничивает его домен, который может быть дополнительно ограничен, если атрибут категоричен.
Например, в наборе данных студентов атрибут может принимать только значения $BS$, $BA$, $MS$ или $PhD$.
Домен атрибута --- $\{BS, BA, MS, PhD\}$.
Атрибут считается категориальным, если во входном наборе данных наблюдается ограниченное количество различных значений для этого атрибута.

Этап предварительной обработки данных продемонстрирован на рисунке~\ref{img:anal-algos}

\includeimage
{anal-algos} % Имя файла без расширения (файл должен быть расположен в директории inc/img/)
{f} % Обтекание (с обтеканием)
{h!} % Положение рисунка (см. wrapfigure из пакета wrapfig)
{0.3\textwidth} % Ширина рисунка
{Этап предварительной обработки данных} % Подпись рисунка
% \imgs{index-struct}{h!}{1}{Индекс как структура данных}
\clearpage

% \section{Особенности метода}

\section{Построение структуры БС с ДП}

На листинге представлен алгоритм, который строит БС $\mathcal{N}$ из входного набора данных $D$, атрибутов $A$ и максимального количества родителей узла БС $k$. 
В этом алгоритме $V$ --- это набор посещаемых атрибутов, а $\Pi$ - это подмножество $V$, которое станет родителями узла $X$, если будет добавлено к $\mathcal{N}$.
Атрибуты $\Pi$ определяются путем поиска максимального $I(X , \Pi)$, который вычисляется по формуле~\ref{eq:dp}:
\begin{equation}\label{eq:dp}
	I(X_i; \Pi) = \sum_{x_i, \pi}^{} P(x_i, \pi) \log \frac{P(x_i, \pi)}{P(x_i) P(\pi)} + Lap(\frac{1}{n\epsilon})
\end{equation}
	, где $\epsilon$ --- параметр приватности, $n$ --- количество записей входных данных.

%На вход подается исходный датасет $D$ с $n$ записями и $d$ атрибутами $A = \{A_1, \dots, A_n\}$, а также параметры:
%\begin{enumerate}[label*=---]
	%\item $k$ --- максимальное число родителей у узла;
	%\item $\epsilon$ --- параметр приватности;
	%\item $categorical\_threshold$ --- порог для категориальных атрибутов.
%\end{enumerate}


%На первом этапе создается пустая сеть и выбирается корневой узел.
%Выбор узла случайный, т.к. порядок не влияет на итоговую структуру благодаря %жадному алгоритму.

%На втором этапе производится рекурсивное добавление узлов.
%Для каждого оставшегося атрибута $X_i \in A \backslash V$:
%\begin{enumerate}
	%item генерация кандидатов в родители для подмножеств $\Pi \subseteq V %\text{размера} \le k$;
	%\item вычисление взаимной информации для каждой пары $(X_i, \Pi) $ происходит по %формуле~\eqref{eq:dp};

	%\item добавляется шум Лапласа, который вычисляется по формуле~\eqref{eq:laplas};
	%\begin{equation}\label{eq:laplas}
	%	\tilde{I}(X_i;\Pi) = I(X_i;\Pi) 
	%\end{equation}

	%$I$ при добавлении/удалении одной записи;
	%\item выбирается лучшая пара $(X_i, \Pi^*)$ с максимальным $\tilde{I}(X_i;\Pi)$, %добавляется узел $(X_i, \Pi^*)$ в $\mathcal{N}$ и обновляется $V = V \cup \{X_i\}$.
% \end{enumerate}

\begin{algorithm}
	\label{alg:preproceses}
	\caption{Preprocess}
	\small
	\SetKwInOut{Input}{Input}
	\SetKwInOut{Output}{Output}
	
	\Input{
		\Indp
		$D$~: датасет;\\
		$A$~: набор атрибутов;\\
		$k$~: максимальное количество родительских узлов;\\
	}
	\Output{
		\Indp
		$N$~: структура Байесовской сети.\\
	}
	
	\Begin{
		$\mathcal{N} \leftarrow \emptyset$ и $V \leftarrow \emptyset$ \\
		$X_1$ $\leftarrow$ $_R A$ \\
		$A \cup  \{(X_1, \emptyset)\}$ и $\mathcal{N} \cup X_1$
		
		\For{$i \gets 2 \text{ to } |A|$}{
			$\Omega \gets \emptyset$ \\
			$p \gets \min(k, |V|)$ \\
			\For{$\text{each } X \in A\backslashV \text{ and each }\Pi \in $\begin{pmatrix}V \\ p\end{pmatrix} }{
							$\Omega \cup \{(X, \Pi)\}$ 
			}
			$(X_i, \Pi_i) = \underset{(X,\Pi) \in \Omega}{\mathrm{argmax}} \ I(X; \Pi)$ \\
			$\mathcal{N} \cup {(X_i, \Pi)}$
			
		}

		\Return $\mathcal{N}$
	}
\end{algorithm}

\newpage
\subsection{Генерация синтетических данных}

На начальном этапе генерации метод читает построенную БС.
Корневые атрибуты, которые не зависят от других переменных, сэмплируются из их безусловных распределений, а дочерние --- из условных.
Обход сети происходит в топологическом порядке.

% Учитывая файл описания набора данных, созданный DataDescriber, Data-Generator сэмплирует синтетические данные из дистрибутивов в этом файле. Алгоритм 2 описывает процесс генерации данных. При вызове в режиме атрибутов, связанных с cor, DataDescriber выбирает значения атрибутов в соответствующем порядке из байесовской сети. Важным свойством дифференциальной конфиденциальности является то, что ее прочность ухудшается при повторных запросах. Чтобы предотвратить утечку частной информации через противников, неоднократно отправляющих запросы на генерацию данных, системный администратор может назначить уникальное случайное семя для каждого человека, которому требуется синтетический набор данных. Для поддержки этого DataGenerator предоставляет функциональность seed для каждого пользователя.


\begin{algorithm}
	\label{alg:preprocess}
	\caption{Preprocess}
	\small
	\SetKwInOut{Input}{Input}
	\SetKwInOut{Output}{Output}
	
\Input{\\
	\Indp   $n$~: количество записей для генерации;\\
	$S$~: метаданные датасета;\\
	$s$~: начальное число.
}
\Output{\\
	\Indp $keys$~: отсортированный массив нормализованных ключей;\\
	$data$~: отсортированный по массиву ключей массив указателей;\\
	$cdf$~~: массив значений функции распределения.\\~
}

\Begin{
	установить начальное значение = s для псевдо-случайной генерации чисел
	Пример корневого атрибута из безусловного распределения
	сформировать оставшиеся атрибуты из условного распределения
	
	\Return Сформированный датасет
}
\end{algorithm}

%\subsection{IDEF0 уровня A1}

%На рисунке представлена диаграмма IDEF0 уровня A1 разрабатываемого метода.
%На вход подается: исходный датасет, на основе которого будут генерироваться синтетические данные, и параметры конфигурации, которые включают в себя $\epsilon$--параметр приватности и количество генерируемых записей. 

